\section{Introduction}
\label{sec:intro}

\subsection{The Fundamental Problem}

Intelligence requires reasoning about the world beyond what is immediately
visible.
A self-driving vehicle must anticipate the scene around a blind corner;
a robotic arm must predict how a grasped object will appear from the wrist
camera after a planned motion; a surgeon must mentally reconstruct
anatomy from a 2-D scan.
In each case, the core operation is identical: \textbf{synthesising a
plausible visual state at an unobserved vantage point}.

This is the problem of \textbf{Novel View Synthesis~(\nvs)}.
Formally, given a set of reference observations
$\{\mathbf{I}_i, \boldsymbol{\pi}_i\}_{i=1}^{n}$ with known camera poses
$\boldsymbol{\pi}_i$, the goal is to render a photo-realistic image
$\hat{\mathbf{I}}$ at an arbitrary target pose $\boldsymbol{\pi}^{*}$,
faithfully preserving 3-D geometry, surface appearance, and lighting.
The problem is inherently ill-posed when $n$ is small or when the target
pose is far from any reference: an infinite family of 3-D scenes
is consistent with any finite set of 2-D projections.
Resolving this ambiguity requires either dense capture (classical IBR) or
a learned prior over the distribution of plausible scenes (modern neural
and generative approaches).

While \nvs{} was studied in the image-based rendering
era~\citep{levoy1996light,mcmillan1995plenoptic,chen1993view,debevec1996modeling},
the past six years have witnessed an extraordinary paradigm shift.
NeRF~\citep{mildenhall2020nerf} demonstrated that a shallow MLP can represent
complex scenes with view-dependent appearance by differentiably volume-rendering
a continuous radiance field, establishing that neural representations could
faithfully encode scene geometry and appearance from posed photographs alone.
3DGS~\citep{kerbl2023gaussian} subsequently showed that an explicit set of
anisotropic Gaussians can render at real-time frame rates, resolving the
speed bottleneck that had limited NeRF's practical applicability.
Most recently, diffusion-based methods~\citep{liu2023zero1to3,poole2023dreamfusion,yu2024viewcrafter}
have moved the frontier toward \emph{generalisable, generative} synthesis
that requires as few as a single image---or even a text description---as input,
fundamentally decoupling scene understanding from dense observation.

\subsection{From NVS to World Modeling}

We argue that this progression reflects a deeper conceptual shift:
\textbf{\nvs{} can be viewed as a key spatial inference capability underlying
Generative World Modeling}.

\paragraph{Scope.}
This survey focuses on the \emph{visual-generative} branch of world modeling, where geometry, viewpoint consistency, and scene simulation are central. It does not aim to comprehensively cover latent dynamics models, model-based reinforcement learning, or planning-centric world model architectures. Within this scope, \nvs{} captures geometric state prediction under observer motion, but does not by itself solve object-centric dynamics, action consequences, or physical causality---capabilities addressed as the frontier in Sections~\ref{sec:world} and~\ref{sec:open}.

\begin{quote}
  \textbf{Working Definition~(World Model).}\quad
  A \emph{Generative World Model} is a learned system that can predict the
  visual (and optionally physical) state of a scene under arbitrary
  conditions---including viewpoint changes, temporal evolution, and agent
  interactions---in a manner that is \emph{geometrically consistent},
  \emph{physically plausible}, and \emph{generalisable} to scenes and conditions
  outside the training distribution.
\end{quote}

Under this definition, \nvs{} is the special case in which the only varying
condition is camera pose.
A full world model must additionally handle temporal dynamics~(D3),
action conditioning~(D4), and physical plausibility~(D5).
Crucially, \emph{spatial consistency enforced by \nvs{} is a necessary
prerequisite}: a world model without geometric grounding is merely a video
predictor---perceptually plausible in distribution but unreliable for tasks
requiring precise spatial reasoning such as robot manipulation or navigation.
This observation motivates the organizing perspective of this survey: the capabilities developed by
the NVS community provide essential geometric building blocks for visual-generative world modeling,
even as full world modeling additionally requires physical causality, object-centric dynamics, and
action-conditioned planning that NVS alone does not address.

Recent systems are actively closing this gap.
Genie~2~\citep{parkerholder2024genie2} generates interactive 3-D environments
from a single image; Cosmos~\citep{agarwal2025cosmos} trains world foundation
models at internet scale; and ViewCrafter~\citep{yu2024viewcrafter} with
StreetCrafter~\citep{yan2024streetcrafter} demonstrate that camera-controlled
video generation is practical at high resolution.
These developments suggest an emerging trend toward a more unified model that
subsumes NVS, video generation, and world simulation into a single
architecture trained at scale.

\subsection{Relationship to Prior Surveys}

Table~\ref{tab:surveys} compares this work with the most closely related
prior surveys.
Tewari~\etal~\citep{tewari2022advances} and
Gao~\etal~\citep{gao2022nerf_survey} focus on \nerf{} representations
without addressing world modelling.
Wu~\etal~\citep{wu2024recentGS} and
Fei~\etal~\citep{fei2024gssurvey} cover 3DGS but not generative methods or 4-D.
Chen~\etal~\citep{chen2024survey_drivingNVS} cover \nvs{} for autonomous driving
but lack a world-modeling lens.
To our knowledge, \textbf{no existing survey jointly covers all three paradigms
(\nerf{}, \tgs{}, generative) through the unifying lens of Generative World
Modeling}, nor proposes a continuous-spectrum taxonomy from static
reconstruction to 4-D interactive generation.
Our work fills this gap by treating the field as a single coherent research
trajectory rather than a collection of independently evolving sub-communities.

\begin{table}[t]
  \centering
  \caption{Comparison with related surveys.
    \textbf{\checkmark}~covered;\; $\sim$~partially;\; --~not covered.}
  \label{tab:surveys}
  \scriptsize
  \setlength{\tabcolsep}{3pt}
  \begin{tabular}{@{}lcccccccc@{}}
    \toprule
    \textbf{Survey} & \textbf{Year}
      & \textbf{NeRF} & \textbf{3DGS} & \textbf{Generative}
      & \textbf{Dynamic} & \textbf{Generalise}
      & \textbf{4D/WM} & \textbf{Unified} \\
    \midrule
    Tewari~\etal~\citep{tewari2022advances}        & 2022 & \yes & \no  & \pt  & \pt & \pt & \no & \no \\
    Gao~\etal~\citep{gao2022nerf_survey}            & 2022 & \yes & \no  & \no  & \pt & \no & \no & \no \\
    Zhu~\etal~\citep{zhu2023survey_ibr}             & 2023 & \pt  & \no  & \no  & \no & \no & \no & \no \\
    Wu~\etal~\citep{wu2024recentGS}                 & 2024 & \no  & \yes & \no  & \pt & \no & \no & \no \\
    Fei~\etal~\citep{fei2024gssurvey}               & 2024 & \no  & \yes & \pt  & \pt & \no & \no & \no \\
    Chen~\etal~\citep{chen2024survey_drivingNVS}    & 2024 & \pt  & \pt  & \no  & \yes& \no & \pt & \no \\
    Zhang~\etal~\citep{zhang2025feedforward_survey} & 2025 & \pt  & \pt  & \no  & \no & \yes& \no & \no \\
    He~\etal~\citep{he2025gs_apps}                  & 2025 & \no  & \yes & \pt  & \pt & \no & \no & \no \\
    \textbf{Ours}                                   & 2026 & \yes & \yes & \yes & \yes& \yes& \yes& \yes\\
    \bottomrule
  \end{tabular}
\end{table}

\subsection{The Unified Spectrum}

We organise all methods along a \emph{conceptual map} defined by three axes
(Figure~\ref{fig:spectrum}):

\begin{enumerate}[leftmargin=*,label=\textbf{Axis~\arabic*.}]
  \item \textbf{Input regime}: dense multi-view $\to$ sparse-view $\to$
        single image $\to$ text/unconditional.
  \item \textbf{Temporal scope}: static scene $\to$ deformable/dynamic $\to$
        4-D world generation.
  \item \textbf{Controllability}: viewpoint only $\to$ viewpoint~+~time $\to$
        viewpoint~+~time~+~action.
\end{enumerate}

This conceptual map is an organizing lens rather than a claim about a single unified research trajectory: methods nearby in the figure need not be optimized for the same task, trained on the same data distribution, or evaluated under comparable protocols. Nevertheless, the map makes explicit which world-modeling capabilities each method addresses and which it leaves open.
Methods near the origin of the spectrum---dense multi-view static reconstruction---are
comparatively mature under controlled benchmarks, with current systems approaching photorealism within minutes of
training.
The open problems cluster at the far end: sparse-input, temporally extended,
action-conditioned generation that generalises to arbitrary scenes.
By mapping every method we survey to a position on this spectrum, we make
explicit which aspects of the world-modeling desiderata each addresses and
which it leaves open.

\begin{figure}[t]
  \centering
  \begin{tikzpicture}[
      method/.style={font=\tiny,align=center,inner sep=2pt},
    ]
    \draw[-{Stealth},thick] (0,0)--(9.2,0)
      node[right,font=\small]{Input: dense $\to$ sparse $\to$ text};
    \draw[-{Stealth},thick] (0,0)--(0,5.4)
      node[above,font=\small]{Temporal: static $\to$ dynamic $\to$ 4D};

    \fill[blue!6,rounded corners=3pt]   (0.1,0.1) rectangle (3.0,2.0);
    \fill[purple!6,rounded corners=3pt] (0.1,2.1) rectangle (4.5,4.0);
    \fill[orange!7,rounded corners=3pt] (3.1,0.1) rectangle (8.9,4.0);
    \fill[red!8,rounded corners=3pt]    (4.6,3.1) rectangle (8.9,5.2);

    \node[font=\tiny\bfseries,text=blue!60]   at (1.5,1.8) {Classical/NeRF};
    \node[font=\tiny\bfseries,text=purple!70] at (2.3,3.8) {Dynamic NeRF/4DGS};
    \node[font=\tiny\bfseries,text=orange!70] at (5.8,0.5) {Generalizable/Generative NVS};
    \node[font=\tiny\bfseries,text=red!70]    at (7.0,5.0) {4D World Models};

    \node[method] at (0.9,0.9)  {NeRF};
    \node[method] at (2.0,1.4)  {3DGS};
    \node[method] at (1.5,0.5)  {Instant-NGP};
    \node[method] at (2.8,1.0)  {Mip-NeRF360};
    \node[method] at (0.7,3.0)  {D-NeRF};
    \node[method] at (2.2,3.4)  {4DGS};
    \node[method] at (1.5,2.5)  {Nerfies};
    \node[method] at (3.5,0.9)  {pixelNeRF\\IBRNet};
    \node[method] at (4.8,1.5)  {pixelSplat\\MVSplat};
    \node[method] at (5.8,2.2)  {DUSt3R\\MASt3R};
    \node[method] at (6.5,1.0)  {Zero-1-to-3\\SyncDreamer};
    \node[method] at (7.8,1.6)  {DreamFusion\\GaussianDreamer};
    \node[method] at (8.2,0.6)  {LRM\\TriplaneGaussian};
    \node[method] at (5.2,3.7)  {ViewCrafter};
    \node[method] at (6.8,4.5)  {Genie2\\Cosmos};
    \node[method] at (8.0,3.8)  {GAIA-1\\UniSim};

    \draw[-{Stealth[length=9pt]},dashed,thick,red!60]
      (1.5,1.5)--(7.5,4.4)
      node[midway,above,sloped,font=\scriptsize\bfseries,red!70]
        {$\to$ Generative World Modeling};
  \end{tikzpicture}
  \caption{A conceptual map of the NVS landscape. The diagonal arrow marks the evolutionary
    trajectory toward 4-D Generative World Modeling.}
  \label{fig:spectrum}
\end{figure}

\subsection{Contributions and Paper Structure}

\textbf{Contributions:}
(1)~A unified continuous-spectrum taxonomy across three axes (input regime,
temporal scope, controllability);
(2)~Cross-paradigm coverage of $>$180 methods with explicit World Modeling
implications in each chapter;
(3)~A related-survey gap analysis (Table~\ref{tab:surveys}) substantiating
the novelty of the World Modeling lens;
(4)~A five-direction open-problems roadmap.

The remainder is organised as follows.
\S\ref{sec:background} establishes background and a unified formulation.
\S\ref{sec:nerf}--\ref{sec:generalisable} survey the three representation
paradigms (NeRF, 3DGS, feed-forward reconstruction).
\S\ref{sec:generative} covers generative \nvs{} with particular depth on
camera-controlled diffusion.
\S\ref{sec:world} examines 4-D generation, neural driving simulators, and
world foundation models.
\S\ref{sec:datasets}--\ref{sec:eval} review datasets and evaluation protocols.
\S\ref{sec:open} proposes a research roadmap.

% ─────────────────────────────────────────────────────────────────────────────
